\section{Hartree--Fock equations and Koopmans' theorem}\label{app:D.4}

The Hartree--Fock equations \cref{eq:3.4-HF} are the minimum condition for the energy of a Slater determinant with respect to its orbitals~\cite{fock1930naherungsmethode}, and Koopmans' theorem~\cite{koopmans1934zuordnung} follows from deleting one orbital from the determinant. For the electron gas both take closed form. The derivation follows Ref.~\cite[Chap.~2]{giuliani2005quantum}.

\subsection{The energy of a determinant}

Take $N_e$ electrons with clamped nuclei. The Hamiltonian is a one-body part, $h = -\nabla^2/2m_e + U_1(\vect{r})$, with $U_1$ the attraction of the nuclei, summed over the electrons, plus the Coulomb pair potential $U_2(\vect{r} - \vect{r}') = e^2/(4\pi|\vect{r} - \vect{r}'|)$ summed over pairs. Let $|\Phi_0\rangle$ be the determinant \cref{eq:2.3-slater} of $N_e$ orthonormal orbitals $\psi_\gamma(x)$, with $x = (\vect{r}, s)$ and $\int\dd x = \sum_s\int\dd^3r$. Its energy was worked out in \cref{sec:2.3}: the one-body operator contributes $\sum_\gamma\langle\gamma|h|\gamma\rangle$, and the two-body operator contributes the direct term minus the exchange term, summed over pairs of occupied orbitals. In position space,
\begin{multline}
  E[\Phi_0] = \sum_\gamma\int\dd x\,\psi^*_\gamma(x)\,h\,\psi_\gamma(x) \\
  + \frac12\sum_{\gamma\gamma'}\int\dd x\,\dd x'\,U_2(\vect{r} - \vect{r}')\Bigl[|\psi_\gamma(x)|^2|\psi_{\gamma'}(x')|^2 - \psi^*_\gamma(x)\psi^*_{\gamma'}(x')\psi_{\gamma'}(x)\psi_\gamma(x')\Bigr] ,
  \label{eq:D.4-energy}
\end{multline}
where both sums run over the occupied orbitals. The terms with $\gamma = \gamma'$ cancel inside the bracket, so no electron interacts with itself.

The density matrix $\varrho(x, x') = \sum_\gamma\psi_\gamma(x)\psi^*_\gamma(x')$ has the electron density on its diagonal, $\nd(\vect{r}) = \sum_s\varrho(x, x)$, and the direct term is the classical electrostatic energy of that density. The exchange term is $-\tfrac12\int\dd x\,\dd x'\,U_2\,|\varrho(x, x')|^2$, so $\varrho$ carries all the dependence of the energy on the occupied orbitals. It is unchanged when the occupied orbitals are mixed among themselves by a unitary matrix $\mathcal{M}$, $\psi_\gamma \to \sum_{\gamma'}\psi_{\gamma'}\mathcal{M}_{\gamma'\gamma}$, because $\mathcal{M}\mathcal{M}^\dagger = 1$ leaves $\varrho$ unchanged. The determinant itself changes only by the phase $\det\mathcal{M}$. The energy is therefore a function of the occupied subspace, not of the individual orbitals.

\subsection{The variational equations}

By the variational principle, $E[\Phi_0] \ge E_0$ for every determinant, where $E_0$ is the exact ground-state energy. The Hartree--Fock determinant is the one that makes $E[\Phi_0]$ stationary, and its lowest stationary point is the Hartree--Fock ground state. The orbitals must stay orthonormal, and we impose this with a Hermitian matrix of Lagrange multipliers $\lambda_{\gamma'\gamma}$, making stationary
\begin{equation*}
  E[\Phi_0] - \sum_{\gamma\gamma'}\lambda_{\gamma'\gamma}\Bigl(\int\dd x\,\psi^*_\gamma(x)\psi_{\gamma'}(x) - \delta_{\gamma\gamma'}\Bigr) .
\end{equation*}
A complex function and its conjugate can be varied independently. The derivative with respect to $\psi^*_\gamma(x)$ of the one-body term is $h\psi_\gamma(x)$. In the double sum, $\psi^*_\gamma$ appears once in the slot $\gamma$ and once in the slot $\gamma'$, and since $U_2$ is symmetric the two contributions are equal and cancel the factor $\tfrac12$. The direct term gives $\int\dd x'\,U_2(\vect{r} - \vect{r}')\,\nd(\vect{r}')\,\psi_\gamma(x)$, and the exchange term gives $-\int\dd x'\,U_2(\vect{r} - \vect{r}')\,\varrho(x, x')\,\psi_\gamma(x')$. Setting the total derivative to zero,
\begin{equation*}
  h_{\mathrm{HF}}\,\psi_\gamma = \sum_{\gamma'}\lambda_{\gamma'\gamma}\,\psi_{\gamma'} , \qquad
  h_{\mathrm{HF}}\,\psi(x) = \Bigl[-\frac{\nabla^2}{2m_e} + U_1(\vect{r}) + v_{\mathrm{H}}(\vect{r})\Bigr]\psi(x) + \int\dd x'\,\Sigma_{\mathrm{x}}(x, x')\,\psi(x') ,
\end{equation*}
with the Hartree potential and the exchange (Fock) term
\begin{equation*}
  v_{\mathrm{H}}(\vect{r}) = \int\dd^3r'\,\frac{e^2\,\nd(\vect{r}')}{4\pi|\vect{r} - \vect{r}'|}, \qquad
  \Sigma_{\mathrm{x}}(x, x') = -\frac{e^2\,\varrho(x, x')}{4\pi|\vect{r} - \vect{r}'|} .
\end{equation*}
The operator $h_{\mathrm{HF}}$ is Hermitian, since $\Sigma_{\mathrm{x}}(x', x)^* = \Sigma_{\mathrm{x}}(x, x')$. Projecting the equation onto $\psi_{\gamma''}$ gives $\lambda_{\gamma''\gamma} = \langle\gamma''|h_{\mathrm{HF}}|\gamma\rangle$, a Hermitian matrix. A unitary rotation of the occupied orbitals leaves $h_{\mathrm{HF}}$ unchanged, because it depends on them only through $\nd$ and $\varrho$, and we can choose the rotation that diagonalizes $\lambda$. In that basis the result takes its canonical form,
\begin{equation}
  h_{\mathrm{HF}}\,\psi_\gamma(x) = E_\gamma\,\psi_\gamma(x), \qquad
  \nd(\vect{r}) = \sum_s\sum_{\gamma\,\mathrm{occ}}|\psi_\gamma(\vect{r}, s)|^2, \quad
  \varrho(x, x') = \sum_{\gamma\,\mathrm{occ}}\psi_\gamma(x)\psi^*_\gamma(x') .
  \label{eq:D.4-HF}
\end{equation}
The equations are nonlinear, since $h_{\mathrm{HF}}$ is built from its own occupied eigenfunctions, and they are solved by iteration until the orbitals that go in are the orbitals that come out. The occupied set is the $N_e$ eigenfunctions of lowest $E_\gamma$; the other eigenfunctions of the same operator are the empty, or virtual, orbitals.

Acting on an occupied orbital $\psi_\gamma$, the term $\gamma' = \gamma$ of $\Sigma_{\mathrm{x}}$ is $-\int\dd x'\,U_2|\psi_\gamma(x')|^2\psi_\gamma(x)$, which cancels the self-repulsion contained in $v_{\mathrm{H}}$. More generally, $-|\varrho(x, x')|^2/\varrho(x, x)$ is a hole in the density of the other electrons of the same spin around an electron at $x$, the exchange hole. Since $\int\dd x''\,\varrho(x, x'')\varrho(x'', x') = \varrho(x, x')$ for orthonormal orbitals, it integrates over $x'$ to exactly one missing electron. The Fock term $\Sigma_{\mathrm{x}}$ is the first-order self-energy of \cref{sec:4.7} in the Coulomb interaction; \cref{app:D.11} obtains it from a diagram.

\subsection{Orbital energies and Koopmans' theorem}

Taking the expectation value of \cref{eq:D.4-HF} in $\psi_\gamma$ leads to the orbital energy,
\begin{equation*}
  E_\gamma = \langle\gamma|h|\gamma\rangle + \sum_{\gamma'\,\mathrm{occ}}\Bigl[\langle\gamma\gamma'|U_2|\gamma\gamma'\rangle - \langle\gamma\gamma'|U_2|\gamma'\gamma\rangle\Bigr] ,
\end{equation*}
the kinetic and nuclear energy of the orbital plus its direct and exchange interaction with every occupied orbital. Summing $E_\gamma$ over the occupied orbitals counts every pair twice, and comparison with \cref{eq:D.4-energy} gives
\begin{equation}
  E_{\mathrm{HF}} = \sum_{\gamma\,\mathrm{occ}}E_\gamma - \frac12\sum_{\gamma\gamma'\,\mathrm{occ}}\Bigl[\langle\gamma\gamma'|U_2|\gamma\gamma'\rangle - \langle\gamma\gamma'|U_2|\gamma'\gamma\rangle\Bigr] .
  \label{eq:D.4-Etot}
\end{equation}
The total energy is not the sum of the orbital energies, and the difference is the double counting of the interaction.

Now remove the electron in an occupied orbital $a$ and keep all the other orbitals as they are. The new state is a determinant of $N_e - 1$ orbitals, and its energy is \cref{eq:D.4-energy} without every term that contains $a$. These are the one-body term $\langle a|h|a\rangle$ and, in the double sum, the terms with $\gamma = a$ or $\gamma' = a$, which are equal in pairs and together give the full sum over $\gamma'$ of the interaction of $a$ with the others. That is $E_a$. Adding an electron in an empty orbital $r$ adds the same kind of terms, with the sum over the $N_e$ occupied orbitals. With the orbitals frozen, therefore,
\begin{equation}
  E_{N_e-1}^{(a)} - E_{N_e} = -E_a, \qquad E_{N_e+1}^{(r)} - E_{N_e} = E_r .
  \label{eq:D.4-koopmans}
\end{equation}
The occupied eigenvalues are minus the removal energies and the empty ones are the addition energies. This is Koopmans' theorem, and it assumes frozen orbitals. In a real removal the other orbitals relax around the hole, and the new determinant of lowest energy is not the one obtained by deleting a column.

In a crystal the relaxation drops out. The orbitals are Bloch states $\psi_{n\vect{k}}$, each spread over the $N$ cells of the crystal, and removing one electron changes $\nd$ and $\varrho$ in each cell by an amount of relative order $1/N$. The orbitals respond at first order in that change. The energy, however, is stationary with respect to the orbitals, so the energy gained by relaxing them is of second order, $1/N^2$ per orbital, and summed over the orbitals of one band it vanishes as $1/N$. Hartree--Fock band energies are therefore the removal and addition energies of the Hartree--Fock description of the crystal, and their difference across the gap is its fundamental gap.

\subsection{The electron gas}

In jellium, \cref{eq:2.3-jellium}, plane waves solve \cref{eq:D.4-HF} for any occupied set of plane waves, since $\varrho$ is then invariant under translations, and the filled Fermi sphere is the natural set. The density is then uniform, its Hartree potential is cancelled by the background (\cref{we:2-jellium}), and $\varrho(x, x') = \delta_{ss'}\Vol^{-1}\sum_{p < k_{\mathrm{F}}}\ee^{\ii\vect{p}\cdot(\vect{r} - \vect{r}')}$ depends only on $\vect{r} - \vect{r}'$. Acting on a plane wave $\ee^{\ii\vect{k}\cdot\vect{r}}$, the exchange term multiplies it by a number. With $U_2$ written through its transform, $\Vol^{-1}\sum_{\vect{k}'}v_{\mathrm{C}}(k')\,\ee^{-\ii\vect{k}'\cdot\vect{r}}$, and the sum over $\vect{p}$ turned into an integral,
\begin{equation*}
  \Sigma_{\mathrm{x}}(k) = -\frac{1}{\Vol}\sum_{p < k_{\mathrm{F}}}v_{\mathrm{C}}(|\vect{k} - \vect{p}|)
  = -\int_{p < k_{\mathrm{F}}}\frac{\dd^3p}{(2\pi)^3}\,\frac{e^2}{k^2 + p^2 - 2kpu} ,
\end{equation*}
with $u$ the cosine of the angle between $\vect{k}$ and $\vect{p}$. The sum runs over one spin, since $\varrho$ is diagonal in spin: only electrons of the same spin exchange. The angular integral is $\int_{-1}^1\dd u\,(k^2 + p^2 - 2kpu)^{-1} = (kp)^{-1}\ln|(k + p)/(k - p)|$, which leaves
\begin{equation*}
  \Sigma_{\mathrm{x}}(k) = -\frac{e^2}{4\pi^2k}\int_0^{k_{\mathrm{F}}}\dd p\;p\,\ln\Bigl|\frac{k + p}{k - p}\Bigr|
  = -\frac{e^2k_{\mathrm{F}}^2}{4\pi^2k}\int_0^1\dd t\;t\,\ln\Bigl|\frac{x + t}{x - t}\Bigr| , \qquad x = \frac{k}{k_{\mathrm{F}}} .
\end{equation*}
The integrand has the antiderivative $\tfrac12(t^2 - x^2)\ln|(x + t)/(x - t)| + xt$, as differentiation shows: the derivative of the logarithm, $2x/(x^2 - t^2)$, times $\tfrac12(t^2 - x^2)$ is $-x$, which cancels the derivative of $xt$. Between $t = 0$ and $1$ it gives $x + \tfrac12(1 - x^2)\ln|(1 + x)/(1 - x)| = 2xF_{\mathrm{L}}(x)$, with the function $F_{\mathrm{L}}$ of \cref{eq:C.8-FL} evaluated at $x$. Hence
\begin{equation}
  \Sigma_{\mathrm{x}}(k) = -\frac{e^2k_{\mathrm{F}}}{2\pi^2}\,F_{\mathrm{L}}\!\left(\frac{k}{k_{\mathrm{F}}}\right), \qquad
  E^{\mathrm{HF}}_{\vect{k}} = \frac{k^2}{2m_e} + \Sigma_{\mathrm{x}}(k) .
  \label{eq:D.4-gas}
\end{equation}
The exchange energy is most negative at the bottom of the band, $\Sigma_{\mathrm{x}}(0) = -e^2k_{\mathrm{F}}/2\pi^2$, since $F_{\mathrm{L}}(0) = 1$, and it falls to half that value at the Fermi surface, since $F_{\mathrm{L}}(1) = \tfrac12$.

The occupied band is widened, from $E_{\mathrm{F}}$ to $E_{\mathrm{F}} + e^2k_{\mathrm{F}}/4\pi^2 = E_{\mathrm{F}} + \alpha k_{\mathrm{F}}/\pi$. At the valence density of silicon, with $k_{\mathrm{F}} = \SI{3.57}{keV}$ and $E_{\mathrm{F}} = \SI{12.5}{eV}$ (\cref{sec:2.5}), the widening is \SI{8.29}{eV} and the band is \SI{20.8}{eV} wide. And the slope of $F_{\mathrm{L}}$ diverges logarithmically at $x = 1$ (\cref{app:C.8}), so the group velocity $\partial E^{\mathrm{HF}}_{\vect{k}}/\partial k$ is infinite at the Fermi surface and the density of states there vanishes. A sharp Fermi surface turns the long range of the bare exchange, $v_{\mathrm{C}} \propto 1/k^2$, into a singular $\Sigma_{\mathrm{x}}(k)$. In the real metal the other electrons screen the exchange, the logarithm disappears, and the velocity at the Fermi surface stays close to $k_{\mathrm{F}}/m_e$. The screened exchange is the GW self-energy of \cref{sec:4.8}.

The exchange energy of the gas follows from \cref{eq:D.4-Etot}. The kinetic part averages to $\tfrac35E_{\mathrm{F}}$ per electron. Since $\Sigma_{\mathrm{x}}$ is the interaction of one electron with all the others, the interaction energy is half its average over the occupied states: $E_{\mathrm{x}}/N_e = \tfrac12\langle\Sigma_{\mathrm{x}}\rangle$. The average of $F_{\mathrm{L}}(x)$ over the Fermi sphere is $\int_0^1\dd x\,3x^2F_{\mathrm{L}}(x)$. Its first part is $\tfrac12$. The second is $\tfrac34\int_0^1\dd x\,(x - x^3)\ln[(1 + x)/(1 - x)]$; integrating by parts with $-(1 - x^2)^2/4$ as the antiderivative of $x - x^3$, which vanishes at $x = 1$, turns it into $\tfrac34\int_0^1\dd x\,\tfrac12(1 - x^2) = \tfrac14$. The average is $\tfrac34$, and
\begin{equation}
  \frac{E_{\mathrm{x}}}{N_e} = -\frac{3e^2k_{\mathrm{F}}}{16\pi^2} = -\frac{3\alpha k_{\mathrm{F}}}{4\pi} = -\frac{0.916}{r_s}\,\mathrm{Ry} ,
  \label{eq:D.4-Ex}
\end{equation}
with $\mathrm{Ry} = \alpha^2m_e/2 = \SI{13.61}{eV}$ and $k_{\mathrm{F}} = (9\pi/4)^{1/3}/(r_sa_0)$, which follows from $\nd_e = 3/(4\pi r_s^3a_0^3)$ and \cref{eq:2.5-fermi}. At the silicon valence density, $r_s = 2.0$ and $E_{\mathrm{x}}/N_e = \SI{-6.22}{eV}$, comparable in size to the kinetic energy per electron, $\tfrac35E_{\mathrm{F}} = \SI{7.5}{eV}$. The local-density approximation \cref{eq:3.4-LDA} is built on \cref{eq:D.4-Ex}.
